%! Advanced Factoring — Unit 3, Lesson 3.2 — Student Packet Cover Sheet
\documentclass[10pt]{article}
\usepackage{algebra2-article}
\usepackage{algebra2-boxes}
\usepackage{ltablex}
\keepXColumns

\begin{document}

% Full-bleed forest banner
\begin{tikzpicture}[remember picture, overlay]
  \fill[forest] (current page.north west)
    rectangle ([yshift=-0.9in]current page.north east);
\end{tikzpicture}
\vspace{-0.8in}

\begin{center}
  {\color{white}\textbf{\LARGE Algebra 2}}\\[4pt]
  {\color{forestmid}\large\textbf{Unit 3 \quad Polynomial Functions}}\\[6pt]
  {\color{white}\textbf{\large Lesson 3.2 \quad Advanced Factoring}}
\end{center}

\vspace{0.22in}
\namedateperiod
\vspace{0.18in}

\begin{learningtargetbox}
\small
\begin{enumerate}[leftmargin=1.5em, itemsep=5pt]
  \item \ldots{} factor out the \textbf{GCF} first --- in one and two variables, including a negative leading coefficient.
  \item \ldots{} choose a factoring method by the \textbf{number of terms}: two, three, or four.
  \item \ldots{} apply the \textbf{difference of squares}, the \textbf{perfect-square trinomial}, and the new \textbf{sum and difference of cubes} identities.
  \item \ldots{} factor a four-term polynomial by \textbf{grouping}, and recognize \textbf{quadratic form}.
  \item \ldots{} factor \textbf{completely} --- every factor prime over the integers --- and recognize when a polynomial is already prime.
  \item \ldots{} \textbf{check} a factorization by multiplying back, and read a function's \textbf{zeros} off its factored form.
\end{enumerate}
\end{learningtargetbox}

\vspace{0.14in}

\begin{tocbox}
\small
\renewcommand{\arraystretch}{1.5}
\begin{tabularx}{\linewidth}{@{} c l X r @{}}
\rowcolor{forestbg}
\textbf{\#} & \textbf{Component} & \textbf{Description} & \textbf{Score} \\
\midrule
1 & Warm-Up        & Spiral review: GCF, trinomials, difference of squares, and perfect cubes  & \blank{1.2cm} \\
2 & Guided Notes   & GCF, the term-count decision tree, identities, grouping, cubes            & \blank{1.2cm} \\
3 & Group Activity & \emph{Taking Polynomials Apart} --- factoring completely, by tier         & \blank{1.2cm} \\
4 & Exit Ticket    & Quick check: a GCF problem, a grouping problem, and an SOL-style item     & \blank{1.2cm} \\
5 & Homework       & Mixed factoring by term count, plus a cube-identity extension             & \blank{1.2cm} \\
\midrule
  & \multicolumn{2}{r}{\textbf{Total}} & \blank{1.2cm} \\
\end{tabularx}
\end{tocbox}

\vspace{0.10in}

\begin{remindbox}
\small Factoring is just \textbf{multiplication run backward}, so you can always check an answer by
multiplying it out. Two habits make it routine. First, \textbf{take the GCF out before anything
else} --- it shrinks the problem and often uncovers an identity that was hidden. Second, \textbf{count
the terms} and let that pick your tool: \emph{two terms} means a difference of squares $a^2-b^2$ or a
sum/difference of cubes $a^3\pm b^3$; \emph{three terms} means a perfect square or ordinary trinomial
factoring; \emph{four terms} means grouping. Then obey the word that matters most in this lesson:
\textbf{completely}. A correct first step is not a finished answer --- look at every factor again and
ask whether it can still be broken apart. And remember the two odd ones out: $a^2+b^2$ is
\textbf{prime}, but $a^3+b^3$ is \textbf{not}.
\end{remindbox}

\end{document}
